\textbf{Phần 1. Trắc nghiệm lựa chọn}
\setcounterpageref{ex}{0}

\begin{ex} %[3X3V44]
Tính giá trị của biểu thức: $A = 150 \times 4 - 450 : 3$.
\choice
{$400$}
{\True $450$}
{$500$}
{$350$}
\loigiai{
$150 \times 4 = 600$; $450 : 3 = 150$; $600 - 150 = 450$.
\textbf{Đáp án của câu hỏi là B}
}
\end{ex}

\begin{ex} %[3X3V44]
Một cửa hàng có 3 kho thóc, mỗi kho chứa $240\text{ kg}$ thóc. Cửa hàng đã bán được $\dfrac{1}{3}$ tổng số thóc. Cửa hàng còn lại bao nhiêu ki-lô-gam thóc?
\choice
{$240\text{ kg}$}
{$500\text{ kg}$}
{$420\text{ kg}$}
{\True $480\text{ kg}$}
\loigiai{
Tổng thóc: $240 \times 3 = 720\text{ kg}$. Đã bán: $720 : 3 = 240\text{ kg}$. Còn lại: $720 - 240 = 480\text{ (kg)}$.
\textbf{Đáp án của câu hỏi là D}
}
\end{ex}

\begin{ex} %[3X3V44]
Tìm số tự nhiên $x$, biết: $x \times 5 + 125 = 625$.
\choice
{\True $100$}
{$90$}
{$110$}
{$120$}
\loigiai{
$x \times 5 = 625 - 125 = 500 \Rightarrow x = 500 : 5 = 100$.
\textbf{Đáp án của câu hỏi là A}
}
\end{ex}

\begin{ex} %[3X3V44]
Một mảnh vườn hình chữ nhật có chu vi $84\text{ m}$, chiều rộng là $14\text{ m}$. Chiều dài của mảnh vườn đó là
\choice
{$24\text{ m}$}
{$32\text{ m}$}
{\True $28\text{ m}$}
{$26\text{ m}$}
\loigiai{
Nửa chu vi: $84 : 2 = 42\text{ m}$. Chiều dài: $42 - 14 = 28\text{ (m)}$.
\textbf{Đáp án của câu hỏi là C}
}
\end{ex}

\begin{ex} %[3X3V44]
Trong phép chia cho 8, số dư lớn nhất có thể là
\choice
{$8$}
{\True $7$}
{$6$}
{$5$}
\loigiai{
Số dư luôn bé hơn số chia, lớn nhất là $8 - 1 = 7$.
\textbf{Đáp án của câu hỏi là B}
}
\end{ex}

\begin{ex} %[3X3V44]
Số lớn là 90, số bé bằng $\dfrac{1}{5}$ số lớn. Số lớn hơn số bé bao nhiêu đơn vị?
\choice
{$18$}
{$75$}
{$80$}
{\True $72$}
\loigiai{
Số bé: $90 : 5 = 18$. Hiệu: $90 - 18 = 72$ đơn vị.
\textbf{Đáp án của câu hỏi là D}
}
\end{ex}

\begin{ex} %[3X3V44]
Một hình tròn có bán kính $8\text{ cm}$. Đoạn thẳng dài nhất nối hai điểm trên hình tròn đó có độ dài là
\choice
{\True $16\text{ cm}$}
{$8\text{ cm}$}
{$24\text{ cm}$}
{$32\text{ cm}$}
\loigiai{
Đoạn thẳng dài nhất nối hai điểm trên đường tròn chính là đường kính: $8 \times 2 = 16\text{ (cm)}$.
\textbf{Đáp án của câu hỏi là A}
}
\end{ex}

\begin{ex} %[3X3V44]
Có 4 xe tải chở hàng, mỗi xe chở được 115 bao phân bón. Người ta chuyển xuống 180 bao. Trên các xe còn lại bao nhiêu bao phân bón?
\choice
{$260$ bao}
{$300$ bao}
{\True $280$ bao}
{$320$ bao}
\loigiai{
Tổng số bao: $115 \times 4 = 460$ bao. Còn lại: $460 - 180 = 280$ bao.
\textbf{Đáp án của câu hỏi là C}
}
\end{ex}

\begin{ex} %[3X3V44]
Nhiệt độ buổi sáng là $18^\circ\text{C}$, buổi trưa nhiệt độ tăng thêm $8^\circ\text{C}$, buổi tối nhiệt độ lại giảm đi $4^\circ\text{C}$ so với buổi trưa. Nhiệt độ buổi tối là
\choice
{$20^\circ\text{C}$}
{\True $22^\circ\text{C}$}
{$24^\circ\text{C}$}
{$26^\circ\text{C}$}
\loigiai{
Buổi trưa: $18 + 8 = 26^\circ\text{C}$. Buổi tối: $26 - 4 = 22^\circ\text{C}$.
\textbf{Đáp án của câu hỏi là B}
}
\end{ex}

\begin{ex} %[3X3V44]
Tính nhanh giá trị: $125 \times 3 + 125 \times 5$.
\choice
{$875$}
{$1250$}
{$950$}
{\True $1000$}
\loigiai{
$125 \times (3 + 5) = 125 \times 8 = 1000$.
\textbf{Đáp án của câu hỏi là D}
}
\end{ex}

\begin{ex} %[3X3V44]
Cho đoạn thẳng $AB = 72\text{ cm}$. Điểm $M$ là trung điểm $AB$, điểm $N$ là trung điểm $AM$. Độ dài đoạn thẳng $NB$ là
\choice
{\True $54\text{ cm}$}
{$36\text{ cm}$}
{$48\text{ cm}$}
{$60\text{ cm}$}
\loigiai{
$AM = 36\text{ cm}$; $AN = 18\text{ cm}$; $NB = 72 - 18 = 54\text{ (cm)}$.
\textbf{Đáp án của câu hỏi là A}
}
\end{ex}

\begin{ex} %[3X3V44]
Một khối lập phương có chu vi một mặt là $32\text{ cm}$. Tổng độ dài tất cả 12 cạnh của khối lập phương là
\choice
{$64\text{ cm}$}
{$128\text{ cm}$}
{\True $96\text{ cm}$}
{$72\text{ cm}$}
\loigiai{
Cạnh: $32 : 4 = 8\text{ cm}$. Tổng 12 cạnh: $12 \times 8 = 96\text{ (cm)}$.
\textbf{Đáp án của câu hỏi là C}
}
\end{ex}

\textbf{Phần 2. Trắc nghiệm đúng sai}
\setcounterpageref{ex}{0}

\begin{ex} %[3X3V44]
Xét tính đúng sai của các nhận định ôn tập cuối học kì 1:
\choiceTF
{\True Thực hiện phép tính trong ngoặc trước, ngoài ngoặc sau.}
{$1\text{ kg} = 100\text{ g}$.}
{\True Đường kính của hình tròn dài gấp 2 lần bán kính.}
{Một hình tam giác có thể có 2 góc vuông.}
\loigiai{
\begin{itemize}
    \item Ý a) Đúng: Quy tắc cơ bản.
    \item Ý b) Sai: Sai vì $1\text{ kg} = 1000\text{ g}$.
    \item Ý c) Đúng: Tính chất hình tròn.
    \item Ý d) Sai: Sai vì tam giác chỉ có nhiều nhất 1 góc vuông.
\end{itemize}
\begin{itemize}
    \item a) \textbf{Đúng}.
    \item b) \textbf{Sai}.
    \item c) \textbf{Đúng}.
    \item d) \textbf{Sai}.
\end{itemize}
}
\end{ex}

\begin{ex} %[3X3V44]
Một mảnh vườn hình chữ nhật có chiều dài $24\text{ m}$, chiều rộng bằng nửa chiều dài:
\choiceTF
{\True Chiều rộng mảnh vườn là $12\text{ m}$.}
{\True Chu vi mảnh vườn là $72\text{ m}$.}
{Nếu gấp đôi chiều rộng thì mảnh vườn có chu vi không đổi.}
{\True Chia mảnh vườn thành 2 hình vuông cạnh $12\text{ m}$ là hoàn toàn chính xác.}
\loigiai{
\begin{itemize}
    \item Ý a) Đúng: $24 : 2 = 12\text{ m}$.
    \item Ý b) Đúng: $(24 + 12) \times 2 = 72\text{ m}$.
    \item Ý c) Sai: Sai vì kích thước thay đổi thì chu vi tăng.
    \item Ý d) Đúng: $24 : 12 = 2$ hình vuông.
\end{itemize}
\begin{itemize}
    \item a) \textbf{Đúng}.
    \item b) \textbf{Đúng}.
    \item c) \textbf{Sai}.
    \item d) \textbf{Đúng}.
\end{itemize}
}
\end{ex}

\textbf{Phần 3. Trả lời ngắn}
\setcounterpageref{ex}{0}

\begin{ex} %[3X3V44]
Tính giá trị biểu thức: $120 \times 3 - 160 : 4$.
\shortans{320}
\loigiai{
\textbf{Đáp số:} $320$\\
$360 - 40 = 320$.
}
\end{ex}

\begin{ex} %[3X3V44]
Một hình vuông có chu vi $68\text{ cm}$. Tính độ dài cạnh của hình vuông đó theo xăng-ti-mét.
\shortans{17}
\loigiai{
\textbf{Đáp số:} $17$\\
$68 : 4 = 17\text{ (cm)}$.
}
\end{ex}

\begin{ex} %[3X3V44]
Số lớn là 84, số bé là 7. Hỏi số lớn gấp mấy lần số bé?
\shortans{12}
\loigiai{
\textbf{Đáp số:} $12$\\
$84 : 7 = 12$ lần.
}
\end{ex}

\begin{ex} %[3X3V44]
Tính kết quả phép chia: $945 : 9$.
\shortans{105}
\loigiai{
\textbf{Đáp số:} $105$\\
$945 : 9 = 105$.
}
\end{ex}

\begin{ex} %[3X3V44]
Đổi $2\ell - 650\text{ ml}$ ra đơn vị mi-li-lít.
\shortans{1350}
\loigiai{
\textbf{Đáp số:} $1350$\\
$2000 - 650 = 1350\text{ (ml)}$.
}
\end{ex}

\begin{ex} %[3X3V44]
Có 3 thùng bánh, mỗi thùng có 32 hộp bánh. Người ta chia đều số bánh đó cho 4 lớp học sinh. Hỏi mỗi lớp nhận được bao nhiêu hộp bánh?
\shortans{24}
\loigiai{
\textbf{Đáp số:} $24$\\
Tổng: $32 \times 3 = 96$ hộp. Mỗi lớp: $96 : 4 = 24$ hộp.
}
\end{ex}
